%%%%%%%%%%%%%%%%%%%%%%%%%%%%%%%%%%%%%%%%%
% Medium Length Professional CV
% LaTeX Template
% Version 2.0 (8/5/13)
%
% Content adapted from main.tex while retaining the cv_4.tex layout.
%%%%%%%%%%%%%%%%%%%%%%%%%%%%%%%%%%%%%%%%%

\documentclass{resume}

\usepackage[dvipsnames]{xcolor}
\definecolor{TitleColor}{HTML}{17324D}
\usepackage{fontawesome5}
\usepackage{enumitem}
\usepackage{mathpazo}
\usepackage{tabularx}
\usepackage{titlesec}
\usepackage{changepage}
\usepackage{needspace}
\usepackage[
  colorlinks=true,
  urlcolor=TitleColor,
  linkcolor=TitleColor,
  pdfauthor={Shufan Shahi},
  pdftitle={Shufan Shahi - Curriculum Vitae}
]{hyperref}

\usepackage[left=0.5in,top=0.4in,right=0.5in,bottom=0.5in]{geometry}

\setlist[itemize]{
  leftmargin=1.25em,
  itemsep=1pt,
  topsep=2pt,
  parsep=0pt,
  partopsep=0pt
}
\setlength{\parskip}{2pt}
\setlength{\parindent}{0pt}

\newcommand{\tab}[1]{\hspace{.2667\textwidth}\rlap{#1}}
\newcommand{\itab}[1]{\hspace{0em}\rlap{#1}}
\newcommand{\projectheading}[4]{%
  \Needspace{3\baselineskip}% Keep the heading with at least the first detail line.
  \begin{tabularx}{\linewidth}{@{}X@{\hspace{1em}}r@{}}
    \textbf{#1 - #2} $|$ \emph{#3} & {\small #4}
  \end{tabularx}\vspace{-4pt}
}

\name{Shufan Shahi}

% resume.cls prints the third \address first, followed by the first and second.
\address{
  \faPhone\ \href{tel:+8801990935044}{+880 1990-935044} \quad
  \faEnvelope\ \href{mailto:shufanshahi@gmail.com}{shufanshahi@gmail.com}
}
\address{
  \faGithub\ \href{https://github.com/shufanshahi}{github.com/shufanshahi} \quad
  \faLinkedin\ \href{https://www.linkedin.com/in/shahi-shufan-680135289/}{linkedin.com/in/shahi-shufan} \quad
  \faGlobe\ \href{https://shufanshahi.github.io/}{shufanshahi.github.io}
}
\address{\faMapMarker*\ Dhaka, Bangladesh}

% Keep the template's centered header while containing its harmless \centerline warning.
\renewcommand{\printaddress}[1]{%
  \begingroup
    \hfuzz=31pt
    \def\\{\addressSep\ }%
    \centerline{#1}\par
  \endgroup
  \addressskip
}

\makeatletter
\renewcommand{\printname}{%
  \begingroup
    \hfuzz=31pt
    \hfil{\color{TitleColor}\MakeUppercase{\namesize\bfseries\@name}}\hfil
    \nameskip\break
  \endgroup
}
\makeatother

\titleformat{\section}
  {\Large\bfseries\scshape}{}{0em}{\color{TitleColor}}
  [\color{TitleColor}\titlerule]
\titlespacing*{\section}{0pt}{1.1em}{0.65em}

\renewenvironment{rSection}[1]{
  \Needspace{3\baselineskip}% Avoid an orphaned section title without reserving a large block.
  \section{#1}
  \begin{list}{}{
      \setlength{\leftmargin}{0.2em}
  }
  \item[]
}{
  \end{list}
  \vspace{0pt}
}

\begin{document}
\fontsize{11.5pt}{14pt}\selectfont

% \begin{rSection}{Professional Profile}
% Computer Science undergraduate with hands-on experience in data analysis, machine learning, and full-stack development. Built web and mobile applications using React, Node.js, and Firebase. Skilled in relational databases, API integration, and delivering real-world solutions through hackathons and academic projects.
% \end{rSection}

\begin{rSection}{Education}
\textbf{Islamic University of Technology} \hfill \textbf{2022--2026} \\
B.Sc. in Computer Science and Engineering \hfill \textbf{CGPA: 3.58} \\
Coursework: Algorithms, Machine Learning, Deep Learning, Databases, and Networks.

\textbf{Chuadanga Govt. College} \hfill \textbf{2021} \\
Higher Secondary Certificate (HSC) \hfill \textbf{GPA: 5.00/5.00}
\end{rSection}

\begin{rSection}{Experience}
\textbf{Student Researcher} \hfill \textbf{2026--Present} \\
\emph{Elite Research Lab}
\begin{itemize}
  \item Support research projects through literature reviews, experiments, and analysis.
\end{itemize}

\textbf{Intern} \hfill \textbf{2025} \\
\emph{Battery Low Interactive Ltd.}
\begin{itemize}
  \item Learned Flutter development, REST API integration, and applications of fine-tuned ML models.
\end{itemize}

\textbf{Researcher} \hfill \textbf{2025--2026} \\
\emph{IUT Computer Vision Lab}
\begin{itemize}
  \item Researched multimodal AI methods to improve integration and mitigate modality imbalance.
\end{itemize}


\textbf{Senior Executive} \hfill \textbf{2025--2026} \\
\emph{Project Aero}
\begin{itemize}
  \item Contributed to drone R\&D by applying AI-based object tracking for international UAV competitions.
\end{itemize}



\end{rSection}



\Needspace{6\baselineskip}
\begin{rSection}{Projects}

\projectheading{DocuMind}{Document Q\&A}{Python, LangChain, LlamaIndex, ChromaDB, Ollama}{%
  \href{https://github.com/shufanshahi/DocuMind}{\faGithub\ GitHub}%
}
\begin{itemize}
  \item Built a local RAG system for documents with cited answers and hallucination guardrails.
  \item Compared 4 chunking methods and 2 embedding models; added hybrid BM25 and reranking.
\end{itemize}

\projectheading{Benhallu}{Bengali Hallucination Detection}{Python, RAG, BanglaBERT, scikit-learn}{%
  \href{https://github.com/shufanshahi/Benhallu}{\faGithub\ GitHub}%
}
\begin{itemize}
  \item Built a retrieval-augmented pipeline to verify Bengali answers against supplied context or a Bengali QA knowledge base.
  \item Created a hallucination dataset and used fine-tuned BanglaBERT with answer-comparison features to classify responses.
\end{itemize}

\projectheading{Bengali Speaker Diarization}{Speaker Turn Detection}{PyTorch, Silero VAD, ECAPA-TDNN}{%
  \href{https://github.com/shufanshahi/Bengali-Speaker-Diarization}{\faGithub\ GitHub}%
}
\begin{itemize}
  \item Built a pipeline using neural voice activity detection, speaker embeddings, and agglomerative clustering to label timestamped speaker turns.
  \item Estimated speaker counts with silhouette scores and exported diarization results to CSV.
\end{itemize}

\projectheading{Drone Vision}{Aerial Tracking}{PyTorch, YOLOv11X, ByteTrack}{%
  \href{https://github.com/shufanshahi/droneVision}{\faGithub\ GitHub}%
}
\begin{itemize}
  \item Fine-tuned YOLOv11X on VisDrone and used ByteTrack for human counting and multi-object tracking.
\end{itemize}

\projectheading{BD Traffic Monitoring System}{Traffic Analytics}{PyTorch, YOLO11L, Flask, Docker}{%
  \href{https://github.com/shufanshahi/bd-traffic-monitoring-system}{\faGithub\ GitHub}
  \quad
  \href{https://huggingface.co/spaces/shufanshahi/bd-traffic-monitoring-system}{\faRocket\ Live}%
}
\begin{itemize}
  \item Fine-tuned YOLO on Bangladesh traffic data to detect local vehicle types.
\end{itemize}

\projectheading{Biztrack}{Business Analytics}{Next.js, Node.js, PostgreSQL}{%
  \href{https://github.com/shufanshahi/biztrack}{\faGithub\ GitHub}
  \quad
  \href{https://www.youtube.com/watch?v=STNoXqq8kQM}{\faYoutube\ Demo}%
}
\begin{itemize}
  \item Built AI modules for real-time finance tracking, risk alerts, customer segmentation, and forecasting.
\end{itemize}

\projectheading{Postfolio}{Professional Networking}{Spring Boot, Next.js, Docker, RabbitMQ, Stripe}{%
  \href{https://github.com/shufanshahi/Postfolio}{\faGithub\ GitHub}
  \quad
  \href{https://www.youtube.com/watch?v=U6N-ioqvg9c}{\faYoutube\ Demo}%
}
\begin{itemize}
  \item Integrates social networking, portfolio and CV creation, job search, and AI interview preparation.
\end{itemize}

\projectheading{Redrick Routinson}{Campus Automation}{Python, React, Flask, Firebase}{%
  \href{https://github.com/shufanshahi/RedrickRoutinson}{\faGithub\ GitHub}
  \quad
  \href{https://www.youtube.com/watch?v=Fg-GkZrsIzw}{\faYoutube\ Demo}%
}
\begin{itemize}
  \item Built GA scheduling, QR exam attendance, class swapping, events, forums, and lost-and-found.
\end{itemize}
\end{rSection}

\begin{rSection}{Research Experiences}
\Needspace{3\baselineskip}
\textbf{EthMuSAM: Mutual Scoring with SAM for Zero-Shot Anomaly Segmentation} \hfill \textbf{2026}
\begin{itemize}
  \item Combined DINOv2 patch embeddings with SAM-based mutual scoring to localize zero-shot anomalies.
\end{itemize}
\Needspace{3\baselineskip}
\textbf{Prototype-Guided Adaptive Feature Learning for Multimodal Emotion Recognition in Conversation (MERC)} \hfill \textbf{2026}
\begin{itemize}
  \item Used adaptive weighting and modality prototypes to improve multimodal emotion recognition.
\end{itemize}
\Needspace{3\baselineskip}
\textbf{Rethinking Visual Features in Graph-Based Multimodal Emotion Recognition in Conversations (MERC): From Low-Level Pixels to Semantic Representations} \hfill \textbf{2026}
\begin{itemize}
  \item Used semantic visual representations to improve graph-based multimodal emotion recognition.
\end{itemize}

\end{rSection}

\begin{rSection}{Technical Skills}
\textbf{Programming Languages:} Python, C++, Java \\
\textbf{Data and Machine Learning:} PyTorch, TensorFlow, scikit-learn, NumPy, Pandas \\
\textbf{Databases:} MySQL, Firebase, PostgreSQL \\
\textbf{Web Development:} React, Tailwind CSS, Next.js, Spring Boot \\
\textbf{Mobile Development:} React Native, Flutter \\
\textbf{Visualization:} Matplotlib, Seaborn \\
\textbf{Development Tools:} Git, Linux \\
\end{rSection}


\begin{rSection}{Achievements, Awards and Scholarships}
\begin{itemize}[leftmargin=*,label=$\bullet$]
  \item \textbf{5th Place, VAND4 International Competition} --- Zero-Shot Industrial Track \hfill \textbf{2026}
  \item \textbf{First Runner-up, IUT Coderush Hackathon} --- Rapid team application development \hfill \textbf{2025}
  \item \textbf{7th Place, BAIC Datathon} --- Data preprocessing, visualization, and modeling \hfill \textbf{2025}
\end{itemize}
\end{rSection}



\begin{rSection}{Reference}
\textbf{Shahriar Ivan} \\
Assistant Professor, Department of CSE, IUT \\
\textbf{Specialization:} Machine Learning, Deep Learning, Computer Vision \\
\faEnvelope\ \href{mailto:shahriarivan@iut-dhaka.edu}{shahriarivan@iut-dhaka.edu}
\end{rSection}

\end{document}
